\chapter{Introduction}

\section{Motivation and research context}

Earth Observation (EO) has become an important source of information for research, environmental monitoring, regional analysis, and teaching. Satellite missions such as Sentinel-2 provide multispectral observations suitable for land-cover and vegetation analysis, while repeated remote-sensing observations support broader terrestrial research \citep{drusch2012sentinel,mahecha2010terrestrial}. Scientific use of these observations, however, depends on more than access to image files. Researchers must identify products, interpret metadata, manage spatial and temporal constraints, preserve provenance, cope with provider failures, and retain a processing path that can be repeated later.

These requirements create a data-engineering problem. A research group may have access to public EO services but still lack a compact infrastructure that connects provider access, validation, storage, spatial indexing, controlled processing, and reproducible reporting. Large cloud platforms can provide substantial capability, but a university research environment may also need a local path that remains understandable, auditable, and usable when credentials or network access are unavailable. The engineering challenge is therefore not simply to download satellite products. It is to construct a trustworthy path from an external provider response to a research-ready and traceable local representation.

The Kvarken Space Center provides the contextual setting for this thesis. The implementation uses a predefined Kvarken study region and focuses on requirements relevant to a small research and educational environment: standards-based access, reproducibility, transparent failure handling, modest infrastructure requirements, and the ability to demonstrate the processing path without depending on a live production service. The thesis does not evaluate the organizational performance of Kvarken Space Center or student learning outcomes. The case-study framing therefore refers to the institutional and regional context used to motivate and bound the artifact.

\section{Problem statement and research gap}

Standards and public platforms address important parts of EO access. STAC provides a common JSON-based description of spatiotemporal assets, OGC standards support interoperable geospatial interfaces, and the Copernicus Data Space Ecosystem (CDSE) provides access to Copernicus data and cloud-oriented processing services \citep{ogc-api-features,stac-spec,cdse,kovacs2026cdse}. These facilities reduce fragmentation at the provider boundary, but they do not by themselves define how a small research project should validate incoming metadata, preserve raw evidence, bound concurrent work, separate transient from permanent failures, maintain a local spatial catalog, or freeze empirical artifacts for later comparison.

Within the scope of this thesis, the research gap is therefore operational and architectural rather than a claim that equivalent technologies do not exist. The question is how standards-based provider access can be assembled into a small, reproducible, and inspectable research infrastructure with explicit contracts. The system must be usable both with external providers and deterministic offline fixtures. It must avoid hidden retry behavior, unbounded task creation, silent coordinate-system assumptions, and mutable experiment evidence.

\section{Objective and scope}

The objective is to design, implement, and evaluate a hybrid EO data infrastructure that connects external standards-based access with local reproducible storage and processing. In this thesis, \emph{hybrid} refers to that combination of external provider access and local research infrastructure, including an offline execution path based on checked-in fixtures.

The implemented scope includes STAC/CDSE access concepts, typed validation, a provider-neutral \texttt{EOScene} model, OAuth2 token handling, bounded asynchronous ingestion, SHA-256 content addressing, append-only provenance, a SQLite/WAL spatial catalog, metadata quality profiling, a range-based raster/NDVI slice, deterministic experiment reports, and a read-only educational HTTP interface.

The evaluation concerns infrastructure behavior and reproducibility. It does not establish production CDSE performance, full GeoTIFF processing, distributed scalability, scientific EO-product accuracy, or educational learning outcomes.

\section{Research questions}

\begin{enumerate}
  \item How can heterogeneous STAC and CDSE metadata be validated and transformed into a transport-neutral EO scene model while preserving provenance?
  \item How can asynchronous ingestion remain bounded and recover from transient timeouts, OAuth2 expiration, and HTTP 429 responses without retrying malformed payloads?
  \item How effectively can a dependency-free SQLite catalog support regional, temporal, quality, and platform queries under concurrent reader/writer contention?
  \item Which empirical measurements and artifacts are sufficient to make an offline EO pipeline baseline reproducible for research and education?
\end{enumerate}

\section{Contributions}

The thesis makes four related contributions. First, it defines an architectural baseline in which provider-specific behavior is isolated from a typed domain model. Second, it implements an engineering artifact that combines bounded concurrency, provenance, embedded spatial persistence, quality filtering, lightweight raster analysis, and a read-only educational interface. Third, it defines a reproducibility protocol based on frozen fixtures, deterministic reports, versioned source, CI, and dated additive experiment artifacts. Fourth, it makes the limits of the evidence explicit so that synthetic benchmark results are not presented as production-provider performance.

The scientific contribution is therefore architectural and methodological rather than a new remote-sensing algorithm.

\section{Structure of the thesis}

Chapter 2 reviews the EO standards, provider context, provenance, and embedded persistence concepts used by the system. Chapter 3 explains the research methodology and evidence strategy. Chapter 4 presents the architecture and its invariants. Chapter 5 describes implementation details. Chapter 6 reports results. Chapter 7 analyzes the four research questions. Chapter 8 discusses contributions and validity. Chapter 9 concludes the thesis and identifies future work.
